% page 119 of the facsimile (image p-118 of the scan)
% block 160.0 x 246.3 mm ; left margin 42.0 mm
\pgc{20.320mm}{0.000mm}{
\l{}
\l{}
\l{\cel{54}l09.}
\l{}
\l{}
\l{\sou{difuzar}. (\sou{trans}.) Extensar omna-direcione, omna-sinse. -DEFIRS.}
\l{}
\l{\sou{digo}. I. Konstrukturo por preventar inundi ; longa konstrukturo}
\l{\cel{4}destinita a retenar la aquo, sur litoro o sur rivo, alonge}
\l{\cel{4}lageto, an la enireyo di portuo. - II. (\sou{metaf.}) Obstaklo}
\l{\cel{4}opozita a to quo tendencas ekirar de sua limiti. - FIS.}
\l{}
\l{\sou{digestar}. (\sou{trans.}) (\sou{fiziol.})\cel{2}Elaborar (nutrivo) por igar asi-}
\l{\cel{4}milebla la parto nutriva, e separar ica de olta qua esas}
\l{\cel{4}eliminenda. - (\sou{Anke metaf}.) - DEFIS.}
\l{}
\l{\sou{digitalo}. (\sou{bot.}) Planto di qua la floro esas simila a ganto-}
\l{\cel{4}fingro, e la floro uzata por plulentigar la movi di la}
\l{\cel{4}homo-kordio. - L.\cel{1}\sou{digitalis}. - DEFISL.}
\l{}
\l{\sou{digitalino}. Substanco efikera di la digitalo purpura.-DEF.}
\l{}
\l{\sou{digitigrado.}\cel{2}(\sou{zool.}) Qua, pro\cel{2}var la tarso e la metatarso plu}
\l{\cel{4}alta kam normale, apogas su sur sua fingri exkluzive, lor}
\l{\cel{4}marcho. - EFIS.}
\l{}
\l{\sou{digna}. Qua meritas (favoroze o desfavoroze) ulo. - FIS.}
\l{}
\l{\sou{digramo}. (\sou{linguistiko}). Grupo de du literi ; sonuno simpla,}
\l{\cel{4}reprezentata skribe per grupo de du literi. (Kom ex.:}
\l{\cel{4}\sou{sh, ch,}\cel{1}e c.) -DEFIS.}
\l{}
\l{\sou{digresar}.\cel{2}(\sou{netrans}.) (pri\cel{1}\sou{diserto}). Traktar, eskartante su,}
\l{\cel{4}deviacante, de la temo. - (\sou{astron.}) Elongaco maxima di}
\l{\cel{4}planeto inferiora. - EFIS.}
\l{}
\l{\sou{dika}.\cel{2}Konsiderata segun la dimensiono opozata a longeso ed a}
\l{\cel{4}larjeso. - DE}
\l{}
\l{\sou{diklina}. (\sou{bot.}) Di qua la floruli separesas de la florini.}
\l{\cel{4}- DEFIS.}
\l{}
\l{\sou{dikogamo}. (\sou{bot.}) Dicesas pri la flori hermafrodita, en qui}
\l{\cel{4}ne eventas sam-instante la developeso di la stamini e di la}
\l{\cel{4}pistilo.}
\l{}
\l{dikotiledona. (\sou{bot.}) Di qua la embriono havas du lobi (kotile-}
\l{\cel{4}doni). - DEFIS.}
\l{}
\l{dikotoma. (bot.)\cel{2}Di qua la stipi di la rameti e di la pedunkli}
\l{\cel{4}esas dividita - subdividita duope. - DEFIRS.}
\l{}
\l{dikotomio.\cel{2}(\sou{logiko}) Metodo di divido, per subdivido du-unajo.}
\l{\cel{4}- DEFIRS.}
\l{}
\l{\sou{dik}tar. (\sou{trans}.) Dicar (ulo) koram ulu, por ke lu skribez olu}
\l{\cel{4}segun ke la vorti pronuncesas. - DEFIRS.}
}
